\documentclass[11pt,a4paper]{article}
\usepackage[utf8]{inputenc}
\usepackage[spanish]{babel}
\usepackage{amsmath,amssymb}
\usepackage{graphicx}
\usepackage{booktabs}
\usepackage{hyperref}
\usepackage{geometry}
\geometry{top=2.5cm, bottom=2.5cm, left=2.5cm, right=2.5cm}

\title{\textbf{Análisis de Deflexión Elástica en Viga Simplemente Apoyada}}
\author{Alan Rivera}
\date{05 de octubre de 2026}

\begin{document}
\maketitle

\section*{Problema y objetivo}

En el presente informe se plantea resolver un flujo de trabajo a partir de datos iniciales desordenados. El objetivo principal es la creación de un esquema de trabajo reproducible que posea una estructura legible, clara y ordenada.

Además, se evaluará el comportamiento de la deflexión en el centro de la luz de una viga simplemente apoyada sometida a una carga puntual centrada, comparando los datos medidos experimentalmente con las predicciones del modelo elástico lineal de Euler-Bernoulli. \cite{hibbeler2017}.

\section{Datos, Parámetros y Conversiones}
Los parámetros geométricos y mecánicos del caso corresponden a:
\begin{itemize}
    \item Luz entre apoyos ($L$): $4{,}0\,\text{m}$
    \item Base de la sección ($b$): $0{,}2\,\text{m}$
    \item Altura de la sección ($h$): $0{,}4\,\text{m}$
    \item Módulo de elasticidad ($E$): $25\,\text{GPa} = 25 \times 10^6\,\text{kN/m}^2$
\end{itemize}

\subsection{Cálculo del Segundo Momento de Área ($I$)}
Para una sección rectangular, el segundo momento de área respecto al eje neutro se calcula como:
\begin{equation}
I = \frac{b \cdot h^3}{12} = \frac{0{,}2\,\text{m} \cdot (0{,}4\,\text{m})^3}{12} = 1{,}0667 \times 10^{-3}\,\text{m}^4
\end{equation}

\section{Método Teórico y Verificación Dimensional}
La deflexión teórica en el centro de la luz ($\delta_{\text{teo}}$) según la teoría de Euler-Bernoulli viene dada por \cite{timoshenko1975}:
\begin{equation}
\delta_{\text{teo}} = \frac{P L^3}{48 E I}
\end{equation}

\subsection{Verificación de Unidades}
Analizando las dimensiones de la fórmula:
\begin{equation}
[\delta] = \frac{[\text{kN}] \cdot [\text{m}]^3}{\left[\frac{\text{kN}}{\text{m}^2}\right] \cdot [\text{m}]^4} = \frac{[\text{kN}] \cdot [\text{m}]^3}{[\text{kN}] \cdot [\text{m}]^2} = [\text{m}]
\end{equation}
Al multiplicar el resultado por $1000\,\text{mm/m}$, se obtiene la deflexión directamente en milímetros ($\text{mm}$).

\section{Resultados y Discusión}
La Tabla~\ref{tab:resultados} resume la comparación entre los valores medidos y teóricos para cada nivel de carga.

\begin{table}[h!]
\centering
\caption{Comparación de deflexión medida vs. teórica.}
\label{tab:resultados}
\begin{tabular}{ccccc}
\toprule
\textbf{Carga $P$ [kN]} & \textbf{$\delta_{\text{medida}}$ [mm]} & \textbf{$\delta_{\text{teo}}$ [mm]} & \textbf{Diferencia Rel. [\%]} & \textbf{$\delta/P$ [mm/kN]} \\
\midrule
0  & 0,00 & 0,00 & 0,00 & -- \\
5  & 0,26 & 0,25 & 4,00 & 0,0520 \\
10 & 0,50 & 0,50 & 0,00 & 0,0500 \\
15 & 0,76 & 0,75 & 1,33 & 0,0507 \\
20 & 1,01 & 1,00 & 1,00 & 0,0505 \\
25 & 1,28 & 1,25 & 2,40 & 0,0512 \\
30 & 1,54 & 1,50 & 2,67 & 0,0513 \\
35 & 1,77 & 1,75 & 1,14 & 0,0506 \\
40 & 2,05 & 2,00 & 2,50 & 0,0513 \\
\bottomrule
\end{tabular}
\end{table}

\begin{figure}[h!]
\centering
\includegraphics[width=0.85\textwidth]{figures/carga_deflexion.png}
\caption{Relación Carga vs. Deflexión en el centro de la viga.}
\label{fig:grafico}
\end{figure}

Como se observa en la Figura~\ref{fig:grafico}, la respuesta medida presenta una rigidez experimental levemente menor que la teórica, con diferencias relativas que no superan el $4{,}0\%$. La constancia aproximada del cociente $\delta/P \approx 0{,}051\,\text{mm/kN}$ ratifica el comportamiento elástico lineal.

\section{Limitaciones y Conclusiones}
\begin{itemize}
    \item \textbf{Limitación del modelo:} La teoría de Euler-Bernoulli no contempla las deformaciones por cortante (efecto Timoshenko), lo que resulta aceptable en vigas esbeltas ($L/h = 10 \ge 10$), pero puede generar discrepancias si la altura aumentara significativamente.
    \item \textbf{Conclusión:} La hipótesis de linealidad elástica describe adecuadamente los datos observados dentro del rango de cargas analizado ($0$ a $40\,\text{kN}$).
\end{itemize}

\bibliographystyle{plain}
\bibliography{referencias}

\appendix
\section{Declaración de Uso de Inteligencia Artificial Generativa}

\begin{description}
    \item[Herramientas:] Gemini.
    \item[Propósito:] Ayuda con del documento en \LaTeX{} para escribir el codigo y estructuración del repositorio según pauta.
    \item[Tarea realizada:] Automatización del cálculo de la deflexión teórica $\delta_{\text{teo}}$ y generación del código para la figura carga-deflexión.
    \item[Contenido incorporado:] Estructura de tablas, ecuaciones en sintaxis de \LaTeX{} y plantilla para la declaración de IA.
    \item[Verificación:] Verificación visual del gráfico generado.
    \item[Responsabilidad:] Se declara la autoría y comprensión total del material técnico expuesto.
\end{description}

\end{document}